\documentclass[12pt]{article}
\usepackage[utf8]{inputenc}
\usepackage[T1]{fontenc}
\usepackage[margin=1in]{geometry}
\usepackage{setspace}
\usepackage{xcolor}
\usepackage[colorlinks=true, urlcolor=blue!60!black, linkcolor=blue!60!black]{hyperref}
\setlength{\parindent}{0em}
\setlength{\parskip}{0.6em}
\renewcommand{\baselinestretch}{1.0}
\begin{document}
\begin{center}
{\LARGE \textit{Agency for International Development v.\ Alliance for Open Society International, Inc.}, 570 U.S.\ 205 (2013)\par}
\vspace{0.5em}
{\large Case Brief --- Two-Page Edition\par}
\end{center}

\section{Memory Jogger}
Congress may specify the activities it subsidizes, but it cannot compel domestic grant recipients to espouse the Government's viewpoint as their own outside the funded program (\href{https://supreme.justia.com/cases/federal/us/570/205/}{570 U.S. at 214--15, 221}).

\section{Detailed Case Facts}
The Leadership Act of 2003 funded NGOs fighting HIV/AIDS abroad. Besides an unchallenged ban on using funds to promote prostitution (\S~7631(e)), its Policy Requirement (\S~7631(f)) denied funds to any group without a policy explicitly opposing prostitution and sex trafficking. Domestic NGOs such as AOSI wished to remain neutral, fearing they would alienate host governments and the sex workers whose trust HIV prevention requires.

\section{Procedural History}
The Southern District of New York enjoined enforcement; after the agencies issued affiliate guidelines, the injunction was reissued and the Second Circuit affirmed (2--1). The D.C. Circuit had upheld the requirement, creating a split.

\section{Judicial Votes}
Affirmed, 6--2. Roberts wrote for Kennedy, Ginsburg, Breyer, Alito, and Sotomayor; Scalia, joined by Thomas, dissented; Kagan took no part.

\section{Holding}
The Policy Requirement violates the First Amendment as applied to domestic organizations. Conditions that define the limits of a spending program are valid; conditions that leverage funding to regulate speech outside the program are not.

\section{Analysis of Opinions}
\textbf{Majority:} As a direct regulation the requirement would plainly be invalid. Reading \textit{Regan}, \textit{League of Women Voters}, and \textit{Rust} together, Roberts held that because \S~7631(e) already kept federal funds from promoting prostitution, \S~7631(f) went further, ``defining the recipient.'' Affiliate guidelines could not cure it: a separate affiliate cannot voice the recipient's beliefs, and a closely identified one forces ``evident hypocrisy.'' The fungibility and message-confusion defenses also failed. \textbf{Scalia:} the Government may enlist ideological allies to carry out its programs, and there is no coercion, since ``the Government is not forcing anyone to say anything''; any NGO may decline the money. \emph{Later:} \textit{AOSI II} (2020) allowed the requirement to be applied to foreign affiliates, which have no First Amendment rights.

\section{Examples: Future Applications}
\textbf{Same-Side:} overdose-prevention grants conditioned on a clinic's corporate policy opposing drug decriminalization compel belief outside the program. \textbf{Opposite-Side:} a clean-energy program may require its funded curriculum to teach that solar and wind power reduce emissions, leaving grantees free to dissent with private funds. \textbf{Fence-Sitter:} an anti-smoking grant program that disqualifies an NGO accepting tobacco-industry sponsorships in its other divisions.

\section{Critique}
Roberts conceded the program/outside-the-program line is ``hardly clear,'' since Congress can redefine a program to subsume a condition, and the Court never said whether the Government may simply \emph{prefer} applicants who already share its view. Progressive critics add that \textit{AOSI II} later forced U.S. NGOs into the very hypocrisy the Court condemned.

\section{Key Quotations}
\begin{quote}
``In the present context, the relevant distinction that has emerged from our cases is between conditions that define the limits of the government spending program---those that specify the activities Congress wants to subsidize---and conditions that seek to leverage funding to regulate speech outside the contours of the program itself.''
\end{quote}
\noindent --- Chief Justice Roberts, for the Court, 570 U.S. at 214--15.

\end{document}
